\documentclass[11pt]{article}
\usepackage{mathblatt}
% Prüfungsblatt nach mathe-nachhilfe pruefung.md § 5 (erster Bau, 10.10.2026), Hand-Satz; Präambel des Setzers.
% Präambel des Setzers (werkzeuge/setzer.py), wortgleich aus dem Hand-Satz T6B (bau/proben/2026-10-09/kathete-berechnen), dort aus K4W.
\makeatletter
\setlength{\mbpfbe}{0mm}% keine Punkte (Bauregel 6.4)
% eingerückt (Bauregel 4.7, 6.6): \gEin Einzug, \gSchrift eine Stufe kleiner
\newlength{\gEin}\setlength{\gEin}{0pt}
\let\gSchrift\relax
\renewcommand{\pfkreuzzeile}[1]{\par\noindent\hangindent3.2em\hangafter1\hspace*{1.6em}$\square$\ #1\par}
\newcommand{\gkreuz}[1]{$\square$\ #1\hspace{2em}}
% Bauregel 6.7: links, was man ansieht, rechts, was man tut; Spaltengrenze überall an derselben Stelle
\newsavebox{\gbox}\newlength{\gLinks}
\newcommand{\gzwei}[2]{\par\smallskip\noindent\sbox{\gbox}{#1}%
  \setlength{\gLinks}{\dimexpr0.5\textwidth-\gEin-\mbpfnr\relax}%
  \ifdim\wd\gbox>\dimexpr\gLinks-4mm\relax
    \usebox{\gbox}\par\smallskip\noindent\begin{minipage}{\linewidth}#2\end{minipage}%
  \else
    \begin{minipage}[t]{\gLinks}\vspace{0pt}\usebox{\gbox}\end{minipage}%
    \begin{minipage}[t]{\dimexpr\linewidth-\gLinks\relax}\vspace{0pt}\raggedright #2\end{minipage}%
  \fi\par}
\RenewDocumentEnvironment{pfaufg}{m m m}{%
  \begin{lrbox}{\mbaufgabebox}\begin{minipage}[t]{\linewidth}%
  \gSchrift\noindent\hspace*{\gEin}\mbpfrandmarke{#1}%
  \begin{minipage}[t]{\mbpfnr}\strut\textbf{#2}\end{minipage}%
  \begin{minipage}[t]{\dimexpr\linewidth-\gEin-\mbpfnr\relax}\raggedright\strut\ignorespaces}%
 {\end{minipage}%
  \end{minipage}\end{lrbox}%
  \setlength{\mbaufgabehoehe}{\dimexpr\ht\mbaufgabebox+\dp\mbaufgabebox\relax}%
  \par\addvspace{7pt}\Needspace*{\mbaufgabehoehe}%
  \noindent\usebox{\mbaufgabebox}\par\addvspace{7pt}}
\makeatother
% Rechenraum als Karo (Bauregel 6.9): n Rechenschritte = n mal 10 mm
\newcommand{\karo}[1]{\par\smallskip\noindent\begin{tikzpicture}
  \clip (0,0) rectangle ({\linewidth-1mm},{#1*10mm});
  \draw[black!16,line width=0.3pt,step=5mm] (0,0) grid ({\linewidth},{#1*10mm});
  \end{tikzpicture}\par}
\newcommand{\kk}[1]{$\square$\,#1\hspace{0.9em}}
\newcommand{\linie}[1]{\,{\color{black!45}\rule[-1pt]{#1}{0.4pt}}\,}
% Zeile am Ende jedes Durchgangs: Originale klein grau, darüber; dann Vorgänger/Nachfolger (Bauregel 3.4, 6.5)
\newcommand{\blattende}[2]{\par\vfill\noindent\begin{minipage}{\linewidth}{\scriptsize\color{mbgrau}Originale: #1\par}\smallskip
{\footnotesize #2\par}\end{minipage}}
\newcommand{\pkt}{\textperiodcentered{}}

\newcommand{\kennung}{P4Z}

\begin{document}
\pfheftstil
\fancyfoot[C]{\parbox[t]{\textwidth}{\mbfhbox\par\vspace{1.5mm}{\footnotesize\color{mbgrau}{\scriptsize \kennung}\hfill\thepage}}}
\pfheftkopf{Länge mit Pythagoras}{P10}

% ===================== Stufe 1 =====================
\Needspace*{0.45\textheight}\pfstufekopf{Gleichung aufstellen}{in 3 der letzten 5 Prüfungen}

% Bank pythagoras-e1-k2-s6 (Gleichung mit anderen Buchstaben, drei Lagen)
\begin{pfaufg}{}{1.}{}
Schreibe für jedes Dreieck die Gleichung des Satzes des Pythagoras auf. Suche zuerst die Hypotenuse – sie liegt dem rechten Winkel gegenüber.\par\medskip\noindent
\begin{minipage}[b]{0.32\linewidth}\centering
\begin{tikzpicture}[line width=0.7pt,scale=0.85]\coordinate (P1) at (0,0);\coordinate (P2) at (2.6,0);\coordinate (P3) at (0,1.6);\draw (P1) -- (P2) -- (P3) -- cycle;\mbrwmarke{(P1)}{(P2)}{(P3)}
\node[font=\small,anchor=north] at (1.3,-0.08) {$m$};\node[font=\small,anchor=east] at (-0.08,0.8) {$n$};\node[font=\small,anchor=south west] at (1.32,0.82) {$k$};\end{tikzpicture}\par\smallskip
{\small a)}\linie{30mm}
\end{minipage}\hfill
\begin{minipage}[b]{0.32\linewidth}\centering
\begin{tikzpicture}[line width=0.7pt,scale=0.85]\coordinate (P1) at (0,0);\coordinate (P2) at (2.5,0);\coordinate (P3) at (0.9,1.2);\draw (P1) -- (P2) -- (P3) -- cycle;\mbrwmarke{(P3)}{(P1)}{(P2)}
\node[font=\small,anchor=north] at (1.25,-0.08) {$r$};\node[font=\small,anchor=south east] at (0.40,0.62) {$p$};\node[font=\small,anchor=south west] at (1.75,0.62) {$q$};\end{tikzpicture}\par\smallskip
{\small b)}\linie{30mm}
\end{minipage}\hfill
\begin{minipage}[b]{0.32\linewidth}\centering
\begin{tikzpicture}[line width=0.7pt,scale=0.85]\coordinate (P1) at (0,1.25);\coordinate (P2) at (1.8,0);\coordinate (P3) at (2.8,1.44);\draw (P1) -- (P2) -- (P3) -- cycle;\mbrwmarke{(P2)}{(P1)}{(P3)}
\node[font=\small,anchor=north east] at (0.85,0.56) {$s$};\node[font=\small,anchor=north west] at (2.37,0.67) {$t$};\node[font=\small,anchor=south] at (1.40,1.42) {$z$};\end{tikzpicture}\par\smallskip
{\small c)}\linie{30mm}
\end{minipage}
\fusshilfe{\mbox{1:~a) $k^2=m^2+n^2$; b) $r^2=p^2+q^2$; c) $z^2=s^2+t^2$}}
\end{pfaufg}

% 2026-FOR-B1j (EBR-Zwilling 2026-EBR-B1j)
\begin{pfaufg}{nach P10 ’26}{2.}{}
Kreuze an, mit welcher Gleichung du die Länge der Seite $w$ berechnen kannst.
\gzwei{\begin{tikzpicture}[line width=0.7pt,scale=0.85]\coordinate (P1) at (0,0);\coordinate (P2) at (2.8,0);\coordinate (P3) at (2.8,1.7);\draw (P1) -- (P2) -- (P3) -- cycle;\mbrwmarke{(P2)}{(P1)}{(P3)}
\node[font=\small,anchor=north] at (1.4,-0.08) {$u$};\node[font=\small,anchor=west] at (2.88,0.85) {$v$};\node[font=\small,anchor=south east] at (1.36,0.9) {$w$};\end{tikzpicture}}{%
\kk{$w=\sqrt{v^2-u^2}$}\par\smallskip\kk{$w=\sqrt{u^2-v^2}$}\par\smallskip\kk{$w=\sqrt{u^2+v^2}$}}
\fusshilfe{\mbox{2:~$w=\sqrt{u^2+v^2}$}}
\end{pfaufg}

% 2024-OS-B1f
\begin{pfaufg}{nach P10 ’24}{3.}{}
Kreuze an, mit welcher Gleichung du die Länge der Seite $z$ berechnen kannst.
\gzwei{\begin{tikzpicture}[line width=0.7pt,scale=0.85]\coordinate (P1) at (0,0);\coordinate (P2) at (3.0,0);\coordinate (P3) at (0,1.5);\draw (P1) -- (P2) -- (P3) -- cycle;\mbrwmarke{(P1)}{(P2)}{(P3)}
\node[font=\small,anchor=north] at (1.5,-0.08) {$z$};\node[font=\small,anchor=east] at (-0.08,0.75) {$y$};\node[font=\small,anchor=south west] at (1.52,0.77) {$x$};\end{tikzpicture}}{%
\kk{$z=x^2-y^2$}\par\smallskip\kk{$z=\sqrt{x^2+y^2}$}\par\smallskip\kk{$z=x+y$}\par\smallskip\kk{$z=\sqrt{x^2-y^2}$}}
\fusshilfe{\mbox{3:~$z=\sqrt{x^2-y^2}$}}
\end{pfaufg}

% ===================== Stufe 2 =====================
\newpage\pfstufekopf{Hypotenuse gesucht}{in 1 der letzten 5 Prüfungen}

% Bank pythagoras-e1-k2-s1-v1, e1-k2-s2-v1, e1-k2-s3-v1 (Päckchen: Zahl, Lage, krumme Wurzel)
\begin{pfaufg}{}{4.}{}
Berechne die Länge der Hypotenuse.\par\smallskip
\gzwei{\parbox[t]{0.3\linewidth}{\textbf{a)}\ Die Katheten sind $36$~cm und $15$~cm lang.}}{\karo{3}}
\gzwei{\rwbei{C}\dreieck{(0,0)}{(5,0)}{(1.4,2.24)}{45\text{ cm}}{28\text{ cm}}{?}{}{}{}}{%
\textbf{b)}\ Wie lang ist die Seite mit dem Fragezeichen?\karo{3}}
\gzwei{\parbox[t]{0.3\linewidth}{\textbf{c)}\ Die Katheten sind $4$~cm und $9$~cm lang. Runde auf eine Stelle nach dem Komma.}}{\karo{3}}
\fusshilfe{\mbox{4:~a) 39 cm; b) 53 cm; c) 9,8 cm}}
\end{pfaufg}

% 2025-OS-K4a, herausgelöst (nur die Länge; der Winkel steht im Blatt „Seite oder Winkel mit sin, cos, tan“)
\begin{pfaufg}{nach P10 ’25}{5.}{}
Familie Yücel baut an die Stufe vor der Haustür eine feste Rampe für Rollstühle. Berechne die Länge $x$ der Rampe.
\gzwei{\begin{tikzpicture}[line width=0.7pt,scale=0.75]\coordinate (P1) at (0,0);\coordinate (P2) at (5.95,0);\coordinate (P3) at (5.95,0.9);\draw (P1) -- (P2) -- (P3) -- cycle;\mbrwmarke{(P2)}{(P1)}{(P3)}
\node[font=\small,anchor=north] at (2.97,-0.08) {$170$ cm};\node[font=\small,anchor=west] at (6.03,0.45) {$16$ cm};\node[font=\small,anchor=south east] at (2.95,0.48) {$x$};\end{tikzpicture}}{%
\karo{3}}
\fusshilfe{\mbox{5:~170,8 cm}}
\end{pfaufg}

% ===================== Stufe 3 =====================
\Needspace*{0.45\textheight}\pfstufekopf{Kathete gesucht}{in 3 der letzten 5 Prüfungen}

% Bank pythagoras-e2-k3-s1-v1, e2-k3-s2-v1 (Päckchen: glatt, krumm)
\begin{pfaufg}{}{6.}{}
Die Hypotenuse und eine Kathete sind gegeben. Berechne die andere Kathete.\par\smallskip
\gzwei{\parbox[t]{0.3\linewidth}{\textbf{a)}\ Hypotenuse $85$~cm, Kathete $13$~cm}}{\karo{3}}
\gzwei{\parbox[t]{0.3\linewidth}{\textbf{b)}\ Hypotenuse $11$~cm, Kathete $6$~cm. Runde auf eine Stelle nach dem Komma.}}{\karo{3}}
\fusshilfe{\mbox{6:~a) 84 cm; b) 9,2 cm}}
\end{pfaufg}

% 2024-OS-K6a, herausgelöst
\begin{pfaufg}{nach P10 ’24}{7.}{}
Eine Seilbahn führt vom Ort $B$ im Tal zum Ort $A$ auf einem Berg. Der Punkt $F$ liegt senkrecht unter $A$ auf der Höhe von $B$. Die Strecke $\overline{FA}$ ist der Höhenunterschied zwischen $B$ und $A$. Berechne diesen Höhenunterschied.
\gzwei{\begin{tikzpicture}[line width=0.7pt,scale=0.8]\coordinate (A) at (0,2.87);\coordinate (F) at (0,0);\coordinate (B) at (2.55,0);\draw (A) -- (F) -- (B) -- cycle;\mbrwmarke{(F)}{(B)}{(A)}
\node[font=\small,anchor=south] at (0,2.92) {$A$};\node[font=\small,anchor=north east] at (0,0) {$F$};\node[font=\small,anchor=north west] at (2.55,0) {$B$};
\node[font=\small,anchor=north] at (1.27,-0.08) {$255$ m};\node[font=\small,anchor=south west] at (1.32,1.48) {$384$ m};\end{tikzpicture}}{%
\karo{3}}
\fusshilfe{\mbox{7:~287,1 m}}
\end{pfaufg}

% 2022-OS-K5a, herausgelöst (Winkel und Seite a nur für 5b–5d, weggelassen)
\begin{pfaufg}{nach P10 ’22}{8.}{}
Das Dreieck $ABC$ hat keinen rechten Winkel. Die Höhe $h_c$ geht von $C$ senkrecht auf die Seite $\overline{AB}$, ihr Fußpunkt heißt $D$. Weise nach, dass die Strecke $\overline{AD}$ etwa $12{,}0$~m lang ist.
\gzwei{\begin{tikzpicture}[line width=0.7pt,scale=0.36]\coordinate (A) at (0,0);\coordinate (B) at (17.8,0);\coordinate (C) at (12,7.4);\coordinate (D) at (12,0);\draw (A) -- (B) -- (C) -- cycle;\draw[dashed] (C) -- (D);\mbrwmarke{(D)}{(B)}{(C)}
\node[font=\small,anchor=north east] at (A) {$A$};\node[font=\small,anchor=north west] at (B) {$B$};\node[font=\small,anchor=south] at (C) {$C$};\node[font=\small,anchor=north] at (D) {$D$};
\node[font=\small,anchor=south east] at (6,3.75) {$14{,}1$ m};\node[font=\small,anchor=east] at (11.9,3.2) {$h_c = 7{,}4$ m};\end{tikzpicture}}{%
\karo{3}}
\fusshilfe{\mbox{8:~$\sqrt{144{,}05}\approx 12{,}0$ m}}
\end{pfaufg}

% 2026-FOR-K4a, herausgelöst (EBR-Zwilling 2026-EBR-K3a; Winkel ε nur für 4b)
\begin{pfaufg}{nach P10 ’26}{9.}{}
Im Viereck $ABCD$ sind gegenüberliegende Seiten parallel. Die Seite $\overline{AB}$ ist über $B$ hinaus bis $F$ verlängert. Die Höhe $h$ steht bei $F$ senkrecht auf dieser Verlängerung. Berechne die Höhe $h$.
\gzwei{\begin{tikzpicture}[line width=0.7pt,scale=0.075]\coordinate (A) at (0,0);\coordinate (B) at (45,0);\coordinate (F) at (58,0);\coordinate (C) at (58,29.2);\coordinate (D) at (13,29.2);\draw (A) -- (B) -- (C) -- (D) -- cycle;\draw (B) -- (F);\draw[dashed] (F) -- (C);\mbrwmarke{(F)}{(C)}{(B)}
\node[font=\small,anchor=north east] at (A) {$A$};\node[font=\small,anchor=north] at (B) {$B$};\node[font=\small,anchor=north west] at (F) {$F$};\node[font=\small,anchor=south west] at (C) {$C$};\node[font=\small,anchor=south east] at (D) {$D$};
\node[font=\small,anchor=south east] at (51.2,14.6) {$32$ cm};\node[font=\small,anchor=north] at (51.5,-3.2) {$13$ cm};\node[font=\small,anchor=west] at (58.8,14.6) {$h$};\end{tikzpicture}}{%
\karo{3}}
\fusshilfe{\mbox{9:~29,2 cm}}
\end{pfaufg}

% ===================== Stufe 4 =====================
\Needspace*{0.45\textheight}\pfstufekopf{Dreieck versteckt}{in 2 der letzten 5 Prüfungen}

% Bank pythagoras-e3-k2-s1-v1 (Grundfall: Höhe im gleichschenkligen Dreieck)
\begin{pfaufg}{}{10.}{}
Die Figur zeigt ein gleichschenkliges Dreieck. Berechne die Höhe $h$. Die Höhe teilt das Dreieck in zwei rechtwinklige Dreiecke.
\gzwei{\begin{tikzpicture}[line width=0.7pt,scale=0.85]\coordinate (A) at (0,0);\coordinate (B) at (2.3,0);\coordinate (C) at (1.15,4.5);\coordinate (M) at (1.15,0);\draw (A) -- (B) -- (C) -- cycle;\draw[dashed] (C) -- (M);\mbrwmarke{(M)}{(B)}{(C)}
\node[font=\small,anchor=north] at (1.15,-0.08) {$32$ cm};\node[font=\small,anchor=west] at (1.78,2.25) {$65$ cm};\node[font=\small,anchor=east] at (0.52,2.25) {$65$ cm};\node[font=\small,anchor=west] at (1.2,1.6) {$h$};\end{tikzpicture}}{%
\karo{4}}
\fusshilfe{\mbox{10:~63 cm}}
\end{pfaufg}

% Bank pythagoras-e3-k4-s3-v2 (Strecke schräg im Zylinder)
\begin{pfaufg}{}{11.}{}
Ein Messbecher hat die Form eines Zylinders. Innen ist er $14$~cm hoch und hat den Durchmesser $9$~cm. Ein Löffel ist $18$~cm lang. Kann der Löffel ganz im Becher verschwinden? Schau dir den Becher von vorn an: Du siehst ein Rechteck.\par\smallskip
\gzwei{\parbox[t]{0.3\linewidth}{\kk{ja}\kk{nein}}}{\karo{4}}
\fusshilfe{\mbox{11:~nein, längste Strecke 16,6 cm}}
\end{pfaufg}

% 2022-OS-K2c, herausgelöst (Sternchenaufgabe; Becher aus 2a, 2b, 2d weggelassen)
\begin{pfaufg}{nach P10 ’22\,$\star$}{12.}{}
Ein Becher hat die Form eines Zylinders ohne Deckel. Er ist $7$~cm hoch und hat den Radius $3{,}2$~cm. Ein Rührstab steht schräg im Becher: Er berührt unten den Rand des Bodens und lehnt oben am gegenüberliegenden Rand. $2$~cm des Stabs ragen über den Becher hinaus. Bestimme die Länge des Stabs.
\gzwei{\begin{tikzpicture}[line width=0.7pt,scale=0.85]
\draw (0,2.8) ellipse (1.28 and 0.45);\draw (-1.28,0) -- (-1.28,2.8);\draw (1.28,0) -- (1.28,2.8);
\draw (-1.28,0) arc (180:360:1.28 and 0.45);\draw[dashed,line width=0.5pt] (-1.28,0) arc (180:0:1.28 and 0.45);
\draw[line width=1.3pt] (-1.28,0) -- (1.28,2.8) -- ($(-1.28,0)!1.27!(1.28,2.8)$);
\end{tikzpicture}}{%
\karo{5}}
\fusshilfe{\mbox{12:~11,5 cm}}
\end{pfaufg}

\blattende{2026 FOR \pkt{} 1 \pkt{} j \quad 2026 FOR \pkt{} 2 \pkt{} c \quad 2026 FOR \pkt{} 4 \pkt{} a \quad 2025 \pkt{} 2 \pkt{} a \quad 2025 \pkt{} 4 \pkt{} a \quad 2024 \pkt{} 1 \pkt{} f \quad 2024 \pkt{} 6 \pkt{} a \quad 2022 \pkt{} 1 \pkt{} g \quad 2022 \pkt{} 2 \pkt{} c \quad 2022 \pkt{} 5 \pkt{} a \quad 2021 \pkt{} 1 \pkt{} h \quad 2020 \pkt{} 7 \pkt{} a \quad 2019 \pkt{} 2 \pkt{} d \quad 2019 \pkt{} 3 \pkt{} a \quad 2018 \pkt{} 6 \pkt{} d \quad 2017 \pkt{} 1 \pkt{} d \quad 2016 \pkt{} 7 \pkt{} b}{Vorher: Quadrat und Wurzel\quad\textbullet\quad Weiter: Seite oder Winkel mit sin, cos, tan}
\end{document}
